问题二以"上游资源-鱼群"动态轨迹为基础，先构建综合生态状态指数，再用它驱动中层猎物承载力，进而建立长江江豚、长江鲟及入侵物种种群响应模型。我们设置高阻隔、无人工放流和污染三种对照情景。基线情景下，长江江豚与长江鲟种群末值分别为禁渔前的1.133倍和1.140倍；高阻隔情景下分别降至0.901倍和0.856倍；无人工放流时长江鲟降至0.519倍；污染情景下二者进一步降至0.610倍和0.704倍。由这些情景对照可见，江湖通道阻隔会同时抑制两类珍稀物种恢复，人工增殖放流对长江鲟的补偿效应更为显著。